\documentclass[10pt,a4paper,twoside=false,parskip=half,DIV=14]{scrbook}
\DeclareTOCStyleEntry[pagenumberwidth=2em]{chapter}{chapter}

\usepackage[english,brazilian]{babel}
\usepackage{lmodern}
\usepackage{luacode}
\usepackage{amsmath,amssymb,amsthm}
\numberwithin{equation}{section}
\renewcommand\theequation{\arabic{equation}}
\usepackage{mathtools}
\mathtoolsset{showonlyrefs}
\usepackage{xcolor}
\usepackage{thmtools}
\usepackage{hyperref}
\usepackage{graphicx}
\hypersetup{colorlinks=true,linkcolor=blue!50!black,filecolor=magenta!50!black,urlcolor=cyan!50!black}
\usepackage[inline]{enumitem}
\usepackage{multicol}

\providecommand{\IMCLanguage}{en}
\providecommand{\IMCBabelLanguage}{english}
\providecommand{\IMCBookletTitle}{IMC Problem Booklet}
\providecommand{\IMCBookletDescription}{This booklet contains problems from the International Mathematics Competition for University Students.}
\providecommand{\IMCBookletFooterNote}{All problems belong to the International Mathematics Competition for University Students.}
\providecommand{\IMCProblemWord}{Problem}

\pagestyle{plain}
\renewcommand*{\chapterpagestyle}{plain}

\declaretheoremstyle[
  spaceabove=1em plus 4em minus 1em,
  spacebelow=1em plus 4em minus 1em,
  bodyfont=\normalfont,
  headfont=\bfseries,
  notefont=\normalfont\bfseries,
  headpunct=.,
  notebraces={},
]{customproblemstyle}

\declaretheorem[
  style=customproblemstyle,
  name=Problem,
  numbered=no,
]{problem}

\declaretheoremstyle[
  bodyfont=\slshape,
  headfont=\bfseries,
  headpunct=.,
]{remarkstyle}

\declaretheorem[
  style=remarkstyle,
  name=Remark,
  numbered=no,
]{remark}

\begin{luacode*}
local lfs = require("lfs")
imc_language = "en"

function file_exists(name)
    local f = io.open(name, "r")
    if f then
        f:close()
        return true
    end
    return false
end

local function tex_escape(text)
    text = text:gsub("\\", "\\textbackslash{}")
    text = text:gsub("([%%#$&{}_])", "\\%1")
    text = text:gsub("%^", "\\textasciicircum{}")
    text = text:gsub("~", "\\textasciitilde{}")
    return text
end

local function first_existing(path, candidates)
    for _, candidate in ipairs(candidates) do
        local full_path = path .. candidate
        if file_exists(full_path) then
            return full_path
        end
    end
    return nil
end

local function collect_years()
    local years = {}

    for entry in lfs.dir("problems") do
        if entry:match("^%d%d%d%d$") then
            years[#years + 1] = entry
        end
    end

    table.sort(years, function(a, b)
        return tonumber(a) > tonumber(b)
    end)

    return years
end

local function collect_problems(year)
    local problems = {}

    for entry in lfs.dir("problems/" .. year) do
        if entry:match("^%d+$") then
            problems[#problems + 1] = entry
        end
    end

    table.sort(problems, function(a, b)
        return tonumber(a) < tonumber(b)
    end)

    return problems
end

local function problem_dir(year, problem)
    return "problems/" .. year .. "/" .. problem .. "/"
end

local function year_title(year)
    return "IMC " .. year
end

local function problem_file(year, problem)
    local path = problem_dir(year, problem)
    return first_existing(path, {
        "problem_" .. imc_language .. ".tex",
        "problem_en.tex",
        "problem_pt.tex",
        "problem.tex",
    })
end

local function problem_heading(year, problem)
    return problem .. "/" .. year
end

function generate_problems_only()
    for _, year in ipairs(collect_years()) do
        tex.print("\\addchap{" .. tex_escape("IMC " .. year) .. "}")

        for _, problem in ipairs(collect_problems(year)) do
            local file = problem_file(year, problem)
            if file then
                tex.print("\\pagebreak[2]")
                tex.print("\\par\\medskip\\noindent\\textbf{\\IMCProblemWord\\ " .. problem_heading(year, problem) .. ".}")
                tex.print("\\label{problem:" .. year .. "-" .. problem .. "}")
                tex.print("\\input{" .. file .. "}")
                tex.print("\\par\\bigskip")
            end
        end
    end
end
\end{luacode*}

\directlua{imc_language = "\luaescapestring{\IMCLanguage}"}

\newcommand{\IMCRenderBooklet}{%
\selectlanguage{\IMCBabelLanguage}%
\frontmatter
\begin{titlepage}
\thispagestyle{empty}
\centering
\vspace*{0.18\textheight}
{\Huge\bfseries \IMCBookletTitle\par}
\vspace{1.5em}
{\large \IMCBookletDescription\par}
\vspace{1.5em}
{\large \url{https://www.imc-math.org.uk/}\par}
\vfill
{\small \IMCBookletFooterNote\par}
\end{titlepage}

\cleardoublepage
\chapter*{\contentsname}
\thispagestyle{plain}
\markboth{\contentsname}{\contentsname}
\setlength{\columnsep}{1.25cm}
\begin{multicols}{2}
\csname @starttoc\endcsname{toc}
\end{multicols}
\mainmatter
\directlua{generate_problems_only()}
}